% Conference poster from a finished paper (alterlab-paper-2-web), keyless: pdflatex poster.tex
% Replace every <<...>> with the paper's own content (copied, never invented). Size: set
% paperwidth/paperheight below (default 48 x 36 in landscape; e.g. 60in x 40in, or A0
% 841mm x 1189mm portrait). Figures go in figures/; missing files show a framed placeholder.
\documentclass[25pt, margin=0mm, innermargin=18mm, blockverticalspace=12mm, colspace=14mm]{tikzposter}
\geometry{paperwidth=48in, paperheight=36in}
% tikzposter sizes its content area at class load; recompute it for the custom paper size
\makeatletter
\setlength{\TP@visibletextwidth}{\textwidth}\addtolength{\TP@visibletextwidth}{-2\TP@innermargin}
\setlength{\TP@visibletextheight}{\textheight}\addtolength{\TP@visibletextheight}{-2\TP@innermargin}
\makeatother
\usepackage[T1]{fontenc}
\usepackage{helvet}
\renewcommand{\familydefault}{\sfdefault}
\usepackage{graphicx}
\usepackage{qrcode}

% Okabe-Ito blue / vermillion; swap for the lab's brand colours
\definecolor{brand}{HTML}{0072B2}
\definecolor{accent}{HTML}{D55E00}
\definecolorstyle{paper}{
  \definecolor{colorOne}{named}{brand}\definecolor{colorTwo}{named}{white}\definecolor{colorThree}{named}{accent}
}{
  \colorlet{backgroundcolor}{white}\colorlet{framecolor}{white}
  \colorlet{titlebgcolor}{colorOne}\colorlet{titlefgcolor}{white}
  \colorlet{blocktitlebgcolor}{colorOne}\colorlet{blocktitlefgcolor}{white}
  \colorlet{blockbodybgcolor}{white}\colorlet{blockbodyfgcolor}{black}
}
\usetheme{Default}\usecolorstyle{paper}\useblockstyle{Default}
\tikzposterlatexaffectionproofoff

\newcommand{\posterfig}[2]{% {file}{width as fraction of \linewidth}
  \IfFileExists{#1}{\includegraphics[width=#2\linewidth]{#1}}%
  {\fbox{\parbox[c][0.25\textheight][c]{#2\linewidth}{\centering\Large missing: \detokenize{#1}}}}}

\title{<<PAPER TITLE>>}
\author{<<Author One>>\textsuperscript{1}, <<Author Two>>\textsuperscript{2}}
\institute{\textsuperscript{1}<<Affiliation One>> \quad \textsuperscript{2}<<Affiliation Two>>}

\begin{document}
\maketitle[width=0.98\paperwidth]
% QR code and logo in the title band (top right / top left)
\node[anchor=north east, fill=white, inner sep=6mm] at (0.5\paperwidth-25mm, 0.5\paperheight-14mm)
  {\qrcode[height=55mm]{<<https://doi.org/10.xxxx/paper-or-project-url>>}};
\node[anchor=north west] at (-0.5\paperwidth+25mm, 0.5\paperheight-14mm)
  {\IfFileExists{figures/logo.png}{\includegraphics[height=55mm]{figures/logo.png}}{}};

\begin{columns}
  \column{0.30}
  \block{Motivation}{
    <<Two or three sentences: the problem and why it matters.>>
    \vspace{1em}
    \begin{itemize}
      \item <<Gap in prior work>>
      \item <<What this paper does>>
    \end{itemize}
  }
  \block{Method}{
    \begin{center}\posterfig{figures/method.png}{0.95}\end{center}
    <<One-paragraph method summary in plain language.>>
  }

  \column{0.40}
  \block{Key results}{
    \begin{center}\posterfig{figures/results.png}{0.95}\end{center}
    \begin{itemize}
      \item {\color{accent}\textbf{<<number>>}} <<what it measures, vs which baseline>>
      \item {\color{accent}\textbf{<<number>>}} <<what it measures>>
    \end{itemize}
  }

  \column{0.30}
  \block{Take-home}{
    \Large <<One sentence the reader should remember.>>
  }
  \block{Limitations}{
    <<Honest limitations, as stated in the paper.>>
  }
  \block{References}{
    \small
    [1] <<Reference as in the paper's bibliography.>>\\
    Paper: <<DOI or arXiv URL>> \quad Code: <<repository URL>>
  }
\end{columns}
\end{document}
